\chapter{Conclusion}

Physical security for key-holders is not toughness. It
is deciding, in daylight, what one person is allowed to
move, what follows them around, and who you will call if
that person is no longer free to speak.

The less any one human can authorize, the less useful
they are to grab. The less of your life is posted, the
harder you are to find. The more delay and extra people
you put on the reserve, the less an attacker gains by
staying violent.

None of this is a checklist you tick in January. Limits
drift. New signers join. Someone starts tweeting numbers
again. Re-read the threat when the org changes, when
someone becomes famous, when a conference week is coming.

If something already happened, do not reuse the
compromised path because it is familiar. Look after the
people first. Then rebuild from material that was never
in that room.

You will not get this perfect. The job is to make
coercion expensive, recovery possible, and panic
unnecessary. Keep the authority small. Keep the humans
alive. Move through the world without advertising the
keys.
